\documentclass[12pt, letterpaper]{article}

%-------Packages---------
\usepackage{amssymb,amsfonts}
\usepackage{mathptmx}
\usepackage{amsthm}
\usepackage{graphicx}
\usepackage[all,arc]{xy}
\usepackage[colorlinks=true, allcolors=blue]{hyperref}
\usepackage{enumerate}
\usepackage{bbm}
\usepackage{dsfont}
\usepackage{mathrsfs}
\usepackage{tikz-cd}
\usepackage{setspace}
\usepackage{import}
\usepackage[utf8]{inputenc}
\usepackage{hyperref}
\usepackage{tikz}
\usepackage{float}
\usepackage{mdframed}
\usepackage{fancyhdr}
\usepackage[nointegrals]{wasysym}

\usepackage[left=1in,top=1in,right=1in,bottom=1in]{geometry}

\usepackage{mathtools}
\usepackage{setspace}
\usepackage{cleveref}
\usepackage{multicol}

%--------Theorem Environments--------
%theoremstyle{plain} --- default
\newmdtheoremenv{theorem}{Theorem}[subsection]
\newmdtheoremenv{corollary}[theorem]{Corollary}
\newtheorem{proposition}[theorem]{Proposition}
\newmdtheoremenv{lemma}[theorem]{Lemma}
\newtheorem{conj}[theorem]{Conjecture}
\newtheorem{quest}[theorem]{Question}

\theoremstyle{definition}
\newmdtheoremenv{definition}[theorem]{Definition}
\newmdtheoremenv{definitions}[theorem]{Definitions}
\newtheorem{example}[theorem]{Example}
\newtheorem{algorithm}[theorem]{Algorithm}
\newtheorem{rmk}[theorem]{Remark}
\newtheorem{exercise}[theorem]{Exercise}

%----------- my commands -------------
\newcommand{\C}{\mathbb{C}}
\newcommand{\R}{\mathbb{R}}
\newcommand{\Q}{\mathbb{Q}}
\newcommand{\Z}{\mathbb{Z}}
\newcommand{\N}{\mathbb{N}}
\renewcommand{\P}{\mathbb{P}}
\newcommand{\contr}{\Rightarrow \Leftarrow}
\newcommand{\HRule}[1]{\rule{\linewidth}{#1}}
\newcommand{\s}{\(\,\)}
\renewcommand{\t}[1]{\text{#1}}
\newcommand{\ang}[1]{\left\langle #1 \right\rangle}

% Meta data
\setstretch{1.2}

\begin{document}

\pagestyle{fancy}
\fancyhf{}
\fancyhead[LOE]{\slshape{\textbf{Vincent Ferrara}}}
\fancyhead[ROE]{\thepage}

\section*{HW 1}

\subsection*{Problem 1}

\subsubsection*{c)}
For $n>3$
\[D_n=\left\{ e,x,x^2,\dots,x^n-1,y,yx,yx^2,\dots,yx^{n-1} \right\}\]
\[x=\dfrac{360^\circ}{n}\]\\
\begin{proof}[\textbf{Claim:} For all \(1\leq j < n\) we have \(yx^j=x^{n-j}y\)]
    \phantom{text}\\
    \textit{Proof.} Base Case: \(m=1 \implies yx=x^{n-1}y\) \checkmark

    Inductive Step: Assume for \(k \geq 1\) the statement is true

    Consider: 
    \begin{align*}
        yx^{k+1}&=yx^kx &\t{Apply indcutive hypothesis on } yx^k\\
                &=x^{n-k}yx \\
                &=x^{n-k}x^{n-1}y\\
                &=x^{n}x^{n-(k+1)}y\\
                &=x^{n-(k+1)}y
    \end{align*}
\end{proof}
This $\implies x^my=yx^{n-m}$ also\\
\phantom{text}\\
Consider: \[
R = \begin{pmatrix}
\cos\left(\frac{360^\circ }{n}\right) & -\sin\left(\frac{360^\circ }{n}\right) \\
\sin\left(\frac{360^\circ }{n}\right) & \cos\left(\frac{360^\circ }{n}\right)
\end{pmatrix}
\]
\[
S = \begin{pmatrix}
1 & 0 \\
0 & -1
\end{pmatrix}
\]
\[
S R^j = \begin{pmatrix}
1 & 0 \\
0 & -1
\end{pmatrix}
\begin{pmatrix}
\cos\left(\frac{360^\circ j}{n}\right) & -\sin\left(\frac{360^\circ j}{n}\right) \\
\sin\left(\frac{360^\circ j}{n}\right) & \cos\left(\frac{360^\circ j}{n}\right)
\end{pmatrix}
= \begin{pmatrix}
\cos\left(\frac{360^\circ j}{n}\right) & -\sin\left(\frac{360^\circ j}{n}\right) \\
-\sin\left(\frac{360^\circ j}{n}\right) & -\cos\left(\frac{360^\circ j}{n}\right)
\end{pmatrix}
\]
Let $0 \leq j < n$
\[
R^{-j} S = \begin{pmatrix}
\cos\left(-\frac{360^\circ j}{n}\right) & -\sin\left(-\frac{360^\circ j}{n}\right) \\
\sin\left(-\frac{360^\circ j}{n}\right) & \cos\left(-\frac{360^\circ j}{n}\right)
\end{pmatrix}
\begin{pmatrix}
1 & 0 \\
0 & -1
\end{pmatrix}
\]
Since \( \cos(-\theta) = \cos(\theta) \) and \( \sin(-\theta) = -\sin(\theta) \), this becomes:
\[
R^{-j} S = \begin{pmatrix}
\cos\left(\frac{360^\circ j}{n}\right) & \sin\left(\frac{360^\circ j}{n}\right) \\
-\sin\left(\frac{360^\circ j}{n}\right) & -\cos\left(\frac{360^\circ j}{n}\right)
\end{pmatrix}
\]
The matrix multiplication gives:
\[
= \begin{pmatrix}
\cos\left(\frac{360^\circ j}{n}\right) & -\sin\left(\frac{360^\circ j}{n}\right) \\
-\sin\left(\frac{360^\circ j}{n}\right) & -\cos\left(\frac{360^\circ j}{n}\right)
\end{pmatrix}
\]
which is the same as the left-hand side.\\
Thus, we have shown that:
\[
S R^j = R^{-j} S
\]
\newpage
To show \( R^{-j} = R^{n-j} \), consider:
Let $0 \leq j < n$
\[
R^{-j} = \begin{pmatrix}
\cos\left(-\frac{360^\circ j}{n}\right) & -\sin\left(-\frac{360^\circ j}{n}\right) \\
\sin\left(-\frac{360^\circ j}{n}\right) & \cos\left(-\frac{360^\circ j}{n}\right)
\end{pmatrix}
\]

\[
R^{n-j} = \begin{pmatrix}
\cos\left(\frac{360^\circ (n-j)}{n}\right) & -\sin\left(\frac{360^\circ (n-j)}{n}\right) \\
\sin\left(\frac{360^\circ (n-j)}{n}\right) & \cos\left(\frac{360^\circ (n-j)}{n}\right)
\end{pmatrix}
\]
Noting that:
\[
\cos\left(\frac{360^\circ (n-j)}{n}\right) = \cos\left(\frac{360^\circ n}{n} - \frac{360^\circ j}{n}\right) = \cos\left(-\frac{360^\circ j}{n}\right)
\]
\[
\sin\left(\frac{360^\circ (n-j)}{n}\right) = \sin\left(\frac{360^\circ n}{n} - \frac{360^\circ j}{n}\right) = -\sin\left(\frac{360^\circ j}{n}\right)
\]
we have:
\[
R^{n-j} = R^{-j}
\]
Subgroup of $GL_2(\R)$ : $\ang{R,S \mid R^n=S^2=I,\; SR=R^{n-1}S}=A$
\begin{itemize}
    \item (Associativity) Given by matrix multiplication
    \item (Identity) \[I = \begin{pmatrix}
        1 & 0 \\
        0 & 1
        \end{pmatrix}\] this is a zero-degree rotation
    \item (Inverses) Each rotation matrix $R^k$ s.t. $1 \leq k < n$ has an inverse $R^{n-k}$. Each
    reflection matrix $SR^j$ s.t. $0 \leq j < n$ has an inverse itself.
    \[R^kR^{n-k}=R^n=I=R^n=R^{n-k}R^k\]
    \[SR^jSR^j=SR^jR^{n-j}S=SIS=SS=I\]
    \item (Closure) \begin{enumerate}
        \item Let $R^j$ and $R^k$ s.t. $1\leq j,k < n$
        \[R^jR^k=\begin{pmatrix}
            \cos\left(\frac{360^\circ j}{n}\right) & -\sin\left(\frac{360^\circ j}{n}\right) \\
            \sin\left(\frac{360^\circ j}{n}\right) & \cos\left(\frac{360^\circ j}{n}\right)
            \end{pmatrix}
            \begin{pmatrix}
                \cos\left(\frac{360^\circ k}{n}\right) & -\sin\left(\frac{360^\circ k}{n}\right) \\
                \sin\left(\frac{360^\circ k}{n}\right) & \cos\left(\frac{360^\circ k}{n}\right)
            \end{pmatrix}
        \]
        Using the angle addition formulas $\cos(a+b) = \cos(a)\cos(b)-\sin(a)\sin(b)$ and $\sin(a+b)=\sin(a)\cos(b)+\cos(a)\sin(b)$
        \[
            = \begin{pmatrix}
                \cos\left(\frac{360^\circ j+k}{n}\right) & -\sin\left(\frac{360^\circ j+k}{n}\right) \\
                \sin\left(\frac{360^\circ j+k}{n}\right) & \cos\left(\frac{360^\circ j+k }{n}\right)
                \end{pmatrix}
        \]
        This is another rotation matrix specifically $R^{j+k \mod n}$
        \item Let $SR^j$ and $R^k$ be a reflection and a rotation s.t. $0 \leq j < n$ and $1 \leq k < n$. Then
        \[(SR^j)R^k=SR^{j+k \mod n}\] This is another reflection in the group.
        \[R^k(SR^j)=(R^kS)R^j=S(R^{n-k}R^j)=SR^{n+j-k \mod n}\] This is another reflection in the group.
        \item Let $SR^j$ and $SR^k$ be reflections s.t. $0 \leq j,k < n$
        \[SR^jSR^k=R^{n-j}SSR^k=R^{n+k-j \mod n}\] This is another rotation in the group.
    \end{enumerate}
\end{itemize}
Let $\varphi : D_n \mapsto A$ s.t. $\varphi(x^j)=R^j$, and $\varphi(yx^j)=SR^j$ s.t. $0 \leq j < n$\\
Well Defined: Le\\
Homomorphism: Let $a,b \in D_n$
\begin{enumerate}
    \item If $a=x^j$ and $b=x^k$ s.t. $1 \leq j,k \leq n$\\
    $\varphi(x^jx^k)=\varphi(x^{j+k})=R^{j+k}=R^jR^k=\varphi(x^j)\varphi(x^k)$
    \item If $a=x^j$ and $b=yx^k$ s.t. $0 \leq j,k < n$\\
    $\varphi(x^jyx^k)=\varphi(yx^{n-j}x^k)=SR^{n+k-j}=R^jSR^{k}=\varphi(x^j)\varphi(yx^k)$\\
    $\varphi(yx^kx^j)=\varphi(yx^{k+j})=SR^{k+j}=SR^kR^{j}=\varphi(yx^k)\varphi(x^j)$\\
    \item If $a=yx^j$ and $b=yx^k$ s.t. $0 \leq j,k < n$\\
    $\varphi(yx^jyx^k)=\varphi(x^{n-j}yyx^k)=\varphi(x^{n+k-j})=R^{n+k-j}=R^{n-j}SSR^{k}=SR^jSR^k=\varphi(yx^j)\varphi(yx^k)$
\end{enumerate}
Injective: Let $a,b \in D_n$ s.t. $\varphi(a)=\varphi(b)$\\
Small Proof Before:
Let $G,H$ be groups and $\varphi$ be a homomorphism from $G$ to $H$ then\\
$\varphi(g)^{-1}=\varphi(g^{-1}) \forall g\in G$
\begin{proof}
    Let $g\in G$\\
    $e_H\varphi(e_G)=e_H\varphi(e_Ge_G)=e_H\varphi(e_G)\varphi(e_G) \implies e_H=\varphi(e_G)$\\
    $e_H=\varphi(e_G)=\varphi(gg^{-1})=\varphi(g)\varphi(g^{-1}) \implies \varphi(g)^{-1}=\varphi(g^{-1})$
\end{proof}
\begin{align*}
    \varphi(a) &= \varphi(b)\\
    \varphi(a)\varphi(b)^{-1} &= I\\
    \varphi(ab^{-1}) &= I\\
\end{align*}
Cases:
\begin{enumerate}
    \item If $a=x^j$ and $b=x^k$ for $0 \leq j,k < n$
    \[
    I=\varphi(x^jx^{-k})=R^jR^{-k} \implies j=k
    \]
    Therefore $a=b$
    \item If $a=yx^j$ and $b=yx^k$ for $0 \leq j,k < n$
    \[
    I=\varphi(yx^jyx^k)=SR^jSR^k=R^{n-j}SSR^k=R^{n+k-j} \implies k=j
    \]
    Therefore $a=b$
    \item If $a=x^j$ and $b=yx^k$ for $0 \leq j,k < n$
    \[
    I=\varphi(x^jyx^k)=R^jSR^k=SR^{n+k-j} \implies R^{j-n-k}=S \contr
    \]
    This is a contradiction since this would turn $A$ into only having $R$
\end{enumerate}
Therefore $a=b$ thus $\varphi$ is Injective\\
Surjective:\\
Let $a\in A$\\
Cases:
\begin{enumerate}
    \item If $a=R^k$ s.t. $0 \leq k < n$
    Consider $\varphi(x^k)=R^k=a$
    \item IF $a=SR^k$ s.t. $0 \leq k < n$
    Consider $\varphi(yx^k)=SR^k=a$
\end{enumerate}
Therefore $\varphi$ is surjective\\
Therefore $\varphi$ is bijective\\
Therefore $\varphi$ is an isomorphism

\newpage
\subsubsection*{d)}
Let $G$ be a group
Consider $Z(G)$
\begin{enumerate}
    \item (Identity)
    Consider $e,g \in G$.\\
    $eg=ge \implies e \in Z(G)$
    \item (Closure)
    Let $z,h \in Z(G)$ and $g \in G$\\
    $zhg=zgh=gzh \implies zh \in Z(G)$
    \item (Inverses)
    Let $z \in Z(G)$ and $g \in G$\\
    $zg=gz \implies g^{-1}z^{-1}=z^{-1}g^{-1} \implies z^{-1}\in Z(G)$
\end{enumerate}
Therefore $Z(G)$ is a subset of $G$

\begin{itemize}
    \item If $n$ is odd then $Z(G)=\{e\}$
    \item If $n$ is even then $Z(G)=\{e,x^{n/2}\}$
\end{itemize}
Proven in c) $x^jy=yx^{n-j}$ for $1 \leq j < n$\\
Thus only $x^j$ commutes with $y$ if $j=n-j \implies j=n/2$\\
and $x^j$ commutes with all $x^k$ since associativity and that $x^kx^j=x^{k+j}$\\
and since $x^j$ commutes with $x^k$ and $y$ it will commute with all reflections\\
Therefore this is only possible if $n$ is even.\\
\phantom{text}\\
Geometrically:\\
When $n$ is odd this is just the do nothing to the $n$-gon.\\
When $n$ is even this is just a $180^\circ$ rotation.

\newpage
\subsubsection*{e)}
Let $G,H$ be a group and $\varphi$ be a homomorphism between them and let $g\in G$
$\forall n \in \Z, \varphi(g^n)=[\varphi(g)]^n$
\begin{proof}
    Base Case: $n=0\;\; \varphi(g^0)=\varphi(e_1)=e_2=[\varphi(g)]^0$

    Assume statement is true for $k\geq0$ for $k \in \N$

    Consider: $k+1$

    $\varphi(g^{k+1})=\varphi(g)\varphi(g^k)=\varphi(g)[\varphi(g)]^k=[\varphi(g)]^{k+1}$

    Note: $e_2=\varphi(e_1)=\varphi(gg^{-1})=\varphi(g)\varphi(g^{-1})\implies [\varphi(g)]^{-1}=\varphi(g^{-1})$

    This proves the statement for all negative integers since we just apply the induction on g then find its inverse
\end{proof}
\phantom{text}\\
If $\varphi$ is an isomorphism then ord$(g)=\text{ord}(\varphi(g))$
\begin{proof}
    Assume $g\in G$ s.t. ord$(g)=n \in N$\\
    $\implies \varphi(g^n)=\varphi(e_G)=e_H \impliedby$ a proof in c)\\
    $\implies \varphi(g)^n=e_H \implies \text{ord}(\varphi(g))$ divides $n$\\
    Since $\varphi$ is injective, if $\varphi(g)^m=e_H$ for some $m<n$ then $g^m=e_G$
    but this contradicts that $g$ has order $n$\\
    Thus ord$(\varphi(g))=$ ord$(g)$
\end{proof}
\phantom{text}\\
Thus, by the proof above an isomorphism has 2 choices
for where $x$ goes either $\tau$ or $\tau^2$ and once you pick $x$ there is only one place for $x^2$\\
and there is 1 choice for $e$\\
and $y$ has 3 choices $\sigma$, $\sigma\tau$, $\sigma\tau^2$, and since
the isomorphism has to preserve group structure there is only one choice for $yx$ and $yx^2$ after you pick $y$\\
Thus, there are in total 6 isomorphisms.\\
They are:
\newpage
\begin{multicols}{3}
\begin{enumerate}
    \item \begin{itemize}
        \item $e$ to $e$
        \item $x$ to $\tau$
        \item $x^2$ to $\tau^2$
        \item $y$ to $\sigma$
        \item $yx$ to $\sigma\tau$
        \item $yx^2$ to $\sigma\tau^2$
        \end{itemize}
    %2
    %3
    \item \begin{itemize}
        \item $e$ to $e$
        \item $x$ to $\tau^2$
        \item $x^2$ to $\tau$
        \item $y$ to $\sigma\tau$
        \item $yx$ to $\sigma$
        \item $yx^2$ to $\sigma\tau^2$
        \end{itemize}
    %5
    \item \begin{itemize}
        \item $e$ to $e$
        \item $x$ to $\tau^2$
        \item $x^2$ to $\tau$
        \item $y$ to $\sigma\tau^2$
        \item $yx$ to $\sigma\tau$
        \item $yx^2$ to $\sigma$
        \end{itemize}
    %7
    \item \begin{itemize}
        \item $e$ to $e$
        \item $x$ to $\tau^2$
        \item $x^2$ to $\tau$
        \item $y$ to $\sigma$
        \item $yx$ to $\sigma\tau^2$
        \item $yx^2$ to $\sigma\tau$
        \end{itemize}
    \item \begin{itemize}
        \item $e$ to $e$
        \item $x$ to $\tau$
        \item $x^2$ to $\tau^2$
        \item $y$ to $\sigma\tau^2$
        \item $yx$ to $\sigma$
        \item $yx^2$ to $\sigma\tau$
        \end{itemize}
    %10
    \item \begin{itemize}
        \item $e$ to $e$
        \item $x$ to $\tau$
        \item $x^2$ to $\tau^2$
        \item $y$ to $\sigma\tau$
        \item $yx$ to $\sigma\tau^2$
        \item $yx^2$ to $\sigma$
        \end{itemize}
    %12
\end{enumerate}
\end{multicols}

Since all of these mappings are injective they are bijective since $|D_3|=|S_3|$.\\
homomorphisms
\begin{multicols}{2}
    \begin{enumerate}
        \item \begin{itemize}
            \item 
        \end{itemize}
    \end{enumerate}
\end{multicols} 


\newpage
\subsection*{Problem 2}
\subsubsection*{a)}

\newpage
\subsection*{Problem 3}
Separate Proof:

        \begin{proof}
            Let $G$ be a finite group and $g\in G$ and let ord$(g)=n\in N$
    
            Consider $\ang{g}$:
    
            \begin{enumerate}
                \item $g^n=e \in \ang{g}$
                \item Let $g_1,g_2 \in \ang{g} \implies g_1=g^j \text{ and } g_2=g^k$ for $j,k\in \N_{\leq n} \text{ thus } g_1g_2=g^jg^k=g^{j+k \mod n} \in \ang{g}$
                \item Let $g_1\in \ang{g} \implies g_1=g^j$ for $j\in \N_{\leq n}$
                
                Consider: $g^{n-j} \in \ang{g}$

                $g^{n-j}g^j=e=g^jg^{n-j} \implies g_1^{-1}=g^{n-j} \in \ang{g}$
            \end{enumerate}
            Thus for all finite groups $G$ and $g\in G$ with finite order, $\ang{g}$ is a subgroup of $G$\\
            $\implies $ by Lagrange's Theorem ord$(g)$ divides ord$(G)$
        \end{proof}
\phantom{text}\\

\begin{proof}
    Let $p$ be prime and $G$ be a group of order $2p$
    \begin{enumerate}
        \item Assume $p=2$\\
        Thus $G$ has 4 elements
        \begin{enumerate}
            \item Assume $G$ has an element with order 4\\
             $\implies \ang{g}$ has $4$ elements since $G$ also has $4$ elements
            and $\ang{g} \subseteq G$ thus there is no element in $G$ that is not in $\ang{g}$ thus $G=\ang{g}$ thus $G \cong C_{4}$
            \item Assume $G$ has no elements with order 4\\
            Thus by my proof above all elements have order 2 for all non-identity elements\\
            Thus $G={e,a,b,c}$ then $a^2=b^2=c^2=e$. WLOG let $c=ab$\\
            Consider $ab=(ab)^{-1}=b^{-1}a^{-1}=ba$\\
            Let $\varphi : G \mapsto D_2$ s.t. $\varphi(e)=e, \varphi(a)=x, \varphi(b)=y, \varphi(ab)=yx=xy=\varphi(ba)$\\
            This is a well-defined function because every input has one output

            Injective:
            
            This is injective by definition.

            Since this is injective and $G$ and $D_2$ have 4 elements this makes $\varphi$ surjective
            Thus it is bijective.\\

            Homomorphism:

            $\varphi(ab)=xy=\varphi(a)\varphi(b)$\\
            $\varphi(ba)=yx=\varphi(b)\varphi(a)$\\
            Let $g\in G$\\
            $\varphi(eg)=\varphi(g)=\varphi(e)\varphi(g)$\\
            $\varphi(ge)=\varphi(g)=\varphi(g)\varphi(e)$\\
            Thus $\varphi$ is an isomorphism thus $G \cong D_2$
        \end{enumerate}
        \item Assume $G$ has an element $g$ s.t. ord$(g)=2p$ for odd primes\\
        $\implies \ang{g}$ has $2p$ elements since $G$ also has $2p$ elements
        and $\ang{g} \subseteq G$ thus there is no element in $G$ that is not in $\ang{g}$ thus $G=\ang{g}$ thus $G \cong C_{2p}$
        \item Assume $G$ has no elements of order $2p$ and $p$ is odd\\
        For contradiction assume $G$ has no elements of order $p \implies$ by my proof above that all non-identity elements have order of 2\\
        Take $e \neq a \neq b \neq e \in G$\\
        Consider $\{e,a,b,ab\} \leq G \implies$ by Lagrange's Theorem that 4 divides $2p \implies 2$ divides $p \contr$\\
        Thus $G$ has an element $g \in G$ with order $p$. Take $\ang{h}=H$ and proven above $H$ is a subgroup of $G$.
        By Lagrange's Theorem $[G:H]=\frac{|G|}{|H|}=\frac{2p}{p}=2$. Take $g\notin H \implies G=H\sqcup gH$
        
        Consider:  $g^2H$
        \begin{enumerate}
            \item Assume $g^2H=H \implies g^2 \in H$\\
            $\implies$ ord$(g^2) \text{ divides } p$ by Lagrange's Theorem $\implies$ ord$(g^2)=1$ or ord$(g)=p$
            \begin{enumerate}
                \item If ord$(g^2)=1 \implies$ $g^{2}=e \implies$ ord$(g)$ divides 2 $\implies$ ord$(g)=2$ or 1
                \begin{enumerate}
                    \item If ord$(g)=1$ then $\ang{g}=\ang{g^2}=\{e\}$ which is a subset of H which $\implies g\in H$ this is a contradiction
                \end{enumerate}
                \item If ord$(g^2)=p \implies$ $g^{2p}=e \implies$ ord$(g)$ divides 2p and since there is no elements of order $2p$\\
                $\implies$ ord$(g)=1$ or $p$
                \begin{enumerate}
                    \item If ord$(g)=p$ then $\ang{g}=\ang{g^2}$ which is a subset of H which $\implies g\in H$ this is a contradiction
                    \item If ord$(g)=1$ then $g=e$ but $g^2\neq e$ since ord$(g^2)=p$ this is a contradiction
                \end{enumerate}
            \end{enumerate}
            Therefore ord$(g)=2$ by 3.(a).i
            \item $g^2H=gH$\\
            Take $h,j \in H$ s.t. $g^2h=gj \implies g=jh^{-1} \in H$ this is a contradiction
        \end{enumerate}
        Therefore anything not in $H$ has order 2.\\
        $gh\in gH \implies gh \notin H \implies$ ord$(gh)=2$\\
        $e=(gh)^2=ghgh \implies h^{-1}g^{-1}=gh \implies gh=h^{p-1}g \implies gh^{k}=h^{p-k}g$
        Since Cosets are a partition\\ 
        $G=H\sqcup gH = \{e,h,h^2,\dots,h^{p-1},g,gh,\dots,gh^{p-1}\}=\ang{h,g \mid e=h^p=g^2, gh=h^{p-1}g}=D_p$
    \end{enumerate}
\end{proof}

\newpage
\subsection*{Problem 4}

\begin{proof}($\implies$)
    Assume $HK$ is a subgroup of $G$\\
    Let $j\in HK$ this means there exists $h\in H$ and $k\in K$ s.t. $j=hk$\\
    Consider $j^{-1}$ since $HK$ is a subgroup of $G$ this implies $j^{-1} \in HK$\\
    Consider $j^{-1}=k^{-1}h^{-1}$ and since $H$ and $K$ are subgroups of $G$ this implies
    that $k^{-1} \in K$ and $h^{-1} \in H$ thus $j^{-1} \in KH$ thus $HK \subseteq KH$\\
    \phantom{text}\\
    Let $j\in KH$ this means there exists $h\in H$ and $k \in K$ s.t. $j=kh$\\
    Consider $j^{-1}=h^{-1}k^{-1}$ and since $H$ and $K$ are subgroups of $G$ this implies
    that $k^{-1} \in K$ and $h^{-1} \in H$ thus $j^{-1} \in HK$ and since $HK$ is a subgroup of $G$
    this implies that $j \in HK$ thus $KH \subseteq HK$\\
    Therefore $HK=KH$
\end{proof}
\begin{proof}($\impliedby$)
    Assume $HK=KH$
    \begin{itemize}
        \item (Identity) Since $H$ and $K$ are subgroups $e \in H\cap K$ thus $ee=e\in HK$
        \item (Inverses) Let $j\in HK \implies \exists h\in H$ and $k\in K$ s.t. $j=hk$\\
        Consider $j^{-1}=k^{-1}h^{-1}$ since $H$ and $K$ are subgroups $k^{-1}\in K$ and $h^{-1}\in H$\\
        therefore $j^{-1}\in KH=HK$
        \item (Closure) Let $j,i \in HK=KH \implies \exists h,h'\in H$ and $k,k'\in K$ s.t. $j=hk$ and $i=h'k'$\\
        $ji=hkh'k'$ since $KH=HK$ there exists an element $n$ in $HK$ s.t. $kh'=n$\\
        $ji=hnk'$  since $n\in HK \implies \exists s \in H$ and $t \in K$ s.t. $n=st$\\
        $ji=hstk'$ since $H$ is a subgroup $hs \in H$ and since $K$ is a subgroup $tk'\in K$ thus $ij\in HK$  
    \end{itemize}
    Therefore $HK$ is a subgroup of $G$.
\end{proof}

\newpage
\subsection*{Problem 5}

   \begin{proof}
     Initially, consider the set of all pairs \( (h, k) \) where \( h \in H \) and \( k \in K \). The total number of such pairs is \( |H| \times |K| \). If all such pairs \( (h, k) \) produced distinct elements in \( HK \), then \( |HK| = |H| \times |K| \). However, since \( H \) and \( K \) may have non-trivial intersection, some products \( hk \) may coincide, leading to overcounting.\\
     Suppose \( h_a k_a = h_b k_b \) for some \( h_a, h_b \in H \) and \( k_a, k_b \in K \). This implies:
     \[
     h_a^{-1} h_b = k_a k_b^{-1}.
     \]
     The left-hand side, \( h_a^{-1} h_b \), is an element of \( H \), and the right-hand side, \( k_a k_b^{-1} \), is an element of \( K \). Thus, \( h_a^{-1} h_b \) must be an element of \( H \cap K \), meaning:
     \[
     h_a^{-1} h_b \in H \cap K \quad \text{and} \quad k_a k_b^{-1} \in H \cap K.
     \]
     Therefore, if \( h_a k_a = h_b k_b \), then the pairs \( (h_a, k_a) \) and \( (h_b, k_b) \) correspond to the same element in \( HK \), and the number of distinct elements in \( HK \) is reduced by a factor equal to the size of the intersection \( H \cap K \).\\
     Therefore, the number of distinct elements in \( HK \) is exactly \( \frac{|H| \cdot |K|}{|H \cap K|} \), establishing that
     \[
     |HK| = \frac{|H| \cdot |K|}{|H \cap K|}.
     \]
   \end{proof}

\end{document}